\documentclass[runningheads]{llncs}

% LNCS basic packages
\usepackage[T1]{fontenc}
\usepackage{graphicx}
\usepackage{amsmath,amssymb}
\usepackage{booktabs}
\usepackage{multirow}
\usepackage{hyperref}
\usepackage{xcolor}
\usepackage{placeins}

% Convenience
\newcommand{\todo}[1]{\textcolor{red}{\textbf{[TODO: #1]}}}

\begin{document}

\title{Topology-aware Low-shot Adaptation of CT Foundation Models\\for Coronary Artery Segmentation}
\titlerunning{Topology-aware Low-shot Adaptation for Coronary Segmentation}

\author{First Author\inst{1}\orcidID{0000-0000-0000-0000} \and
Second Author\inst{1}\orcidID{0000-0000-0000-0000}}
%
\authorrunning{F. Author et al.}
%
\institute{University Name, City, Country \\
    \email{\{first,second\}@university.ac.uk}}

\maketitle

%----------------------------------------------------------------------
\begin{abstract}
Foundation models for CT segmentation achieve strong volumetric overlap through large-scale pretraining, but their adaptation for topology-sensitive tasks such as coronary artery segmentation remains under-studied.
We present a systematic study of how encoder adaptation strategy interacts with topology-aware supervision when fine-tuning CT-FM, a SegResNet foundation model pretrained on 148K CT scans.
On ImageCAS ($n{=}250$ test), we cross five adaptation strategies---decoder-only, BitFit, skip adapters, partial unfreeze, and full fine-tuning---with three loss weights for a differentiable bifurcation connectedness (Soft BCS) surrogate, yielding an ablation grid of segmentation and topology metrics.
Decoder-only fine-tuning achieves the best overlap--connectivity balance (Dice~0.793, BCS~0.716 at $\lambda_s{=}0.03$), while full fine-tuning catastrophically overfits (Dice drops 4.5\,pp, HD95 triples).
Centered Kernel Alignment confirms that all parameter-efficient strategies preserve encoder representations (CKA${\geq}0.84$), localising adaptation to the decoder.
We further evaluate cross-domain transfer to ASOCA, where annotation protocol mismatch (full-tree vs.\ major-vessel labels) causes systematic over-segmentation (Dice~0.48, BCS~0.79).
Minimal low-shot adaptation ($k{=}1$) with shallow decoder fine-tuning and threshold calibration recovers Dice to 0.706 while preserving BCS at 0.752; adding more training data paradoxically degrades performance.
\keywords{Coronary artery segmentation \and Foundation models \and Topology-aware learning \and Low-shot adaptation \and Encoder adaptation}
\end{abstract}

%----------------------------------------------------------------------
\section{Introduction}

Foundation models pretrained on large-scale CT data offer strong initialisation for downstream segmentation tasks, reducing the need for task-specific training from scratch~\cite{pai_ctfm_2025,huang_stu_2023}.
In coronary artery segmentation, accurate delineation underpins stenosis quantification, FFR-CT estimation, and stent planning~\cite{lin_deep_2022}.
However, two practical challenges limit straightforward deployment of pretrained models.

First, \emph{how to adapt} the pretrained encoder is an open question.
Full fine-tuning risks catastrophic forgetting of pretrained features, while freezing the encoder entirely may under-utilise its capacity.
The transfer-learning literature in NLP and 2D vision has studied parameter-efficient strategies---BitFit~\cite{zaken_bitfit_2022}, adapters~\cite{houlsby_adapters_2019}, partial unfreezing~\cite{lee_surgical_2023}---but their interaction with \emph{topology-aware losses} in 3D medical imaging is unexplored.

Second, \emph{annotation protocol mismatch} between source and target datasets creates a distribution shift in label space.
A model pretrained on full coronary tree annotations (ImageCAS~\cite{zeng_imagecas_2023}) will systematically over-segment when applied to datasets labelling only major vessels (ASOCA~\cite{gharleghi_asoca_2022}).
This is not domain shift in image appearance but a mismatch in the definition of the target structure.
In the low-shot regime, where few labelled examples are available, topology-aware supervision may act as an inductive bias to constrain the solution space~\cite{shit_cldice_2021,hu_topology_2019,clough_topological_2022}.

We address both challenges in a unified study.
On ImageCAS ($n{=}638$ train, $n{=}250$ test), we systematically ablate five encoder adaptation strategies crossed with three topology loss weights, using the differentiable Soft BCS loss from our prior work~\cite{anonymous_bcs_2026}.
We complement the ablation with representational analysis (CKA similarity~\cite{kornblith_cka_2019} and linear probing) to explain \emph{why} certain strategies succeed.
We then evaluate cross-domain transfer to ASOCA under annotation protocol mismatch, testing whether the best adaptation strategy generalises and whether topology-aware supervision improves low-shot sample efficiency.

\paragraph{Contributions.}
\begin{enumerate}
    \item A systematic ablation of five encoder adaptation strategies $\times$ three BCS loss weights, identifying decoder-only fine-tuning as optimal for topology-aware coronary artery segmentation.
    \item Representational analysis (CKA, linear probing) showing that all parameter-efficient strategies preserve encoder features (CKA${\geq}0.84$), while full fine-tuning destroys them (CKA${\leq}0.37$ at the bottleneck).
    \item Cross-domain evaluation on ASOCA demonstrating annotation protocol mismatch effects and the role of topology-aware supervision in low-shot adaptation.
\end{enumerate}

%----------------------------------------------------------------------
\section{Related Work}

\paragraph{Foundation model adaptation.}
Large pretrained models can be adapted via full fine-tuning, linear probing, or parameter-efficient methods including BitFit (bias-only tuning)~\cite{zaken_bitfit_2022}, bottleneck adapters~\cite{houlsby_adapters_2019,chen_adaptformer_2022}, and surgical fine-tuning of selected layers~\cite{lee_surgical_2023}.
Kumar et al.~\cite{kumar_finetuning_2022} showed that full fine-tuning can distort pretrained features and degrade out-of-distribution performance.
In medical imaging, STU-Net~\cite{huang_stu_2023} and CT-FM~\cite{pai_ctfm_2025} demonstrate scalable pretraining, but downstream adaptation strategies remain under-explored, particularly for topology-sensitive losses.

\paragraph{Topology-aware segmentation.}
clDice~\cite{shit_cldice_2021} penalises centreline disconnections via soft skeletonisation.
Persistent homology losses~\cite{clough_topological_2022,hu_topology_2019} target Betti numbers but are computationally expensive for 3D volumes.
Roug\'e et al.~\cite{rouge_topology_2024} combine topology-aware cascaded U-Nets for cerebrovascular segmentation.
Our prior work~\cite{anonymous_bcs_2026} introduced bifurcation connectedness score (BCS) and a differentiable Soft BCS loss that directly supervises branch-level connectivity.
All prior topology work focuses on \emph{loss design}; none studies how the choice of \emph{adaptation strategy} interacts with topology-aware objectives.

%----------------------------------------------------------------------
\section{Methods}

\subsection{Background: BCS and Soft BCS Loss}
\label{sec:bcs}

The bifurcation connectedness score (BCS)~\cite{anonymous_bcs_2026} evaluates whether predicted segmentations preserve connectivity at ground-truth bifurcation points.
Given a ground-truth skeleton, bifurcation voxels (degree ${\geq}3$) are clustered into junctions.
For each junction, branch stubs of length $l{=}8$ are traced along each outgoing branch.
A junction is \emph{preserved} if all its stubs are mutually reachable through the predicted segmentation (dilated skeleton, tolerance $\tau{=}3$).
BCS is the fraction of preserved junctions.

The differentiable Soft BCS loss operates on precomputed ground-truth stub labels.
Writing $p(v)$ for the predicted foreground probability at voxel $v$ and $\bar{p}_k = \frac{1}{|s_k|}\sum_{v \in s_k} p(v)$ for the mean probability on stub~$k$, the per-bifurcation loss uses a softmin weighting:
\begin{equation}\label{eq:softbcs}
\ell_c = 1 - \sum_k w_k \bar{p}_k, \quad
w_k = \frac{\exp(-\bar{p}_k/T)}{\sum_j \exp(-\bar{p}_j/T)},
\end{equation}
where $T{=}0.1$ concentrates gradients on the least-connected branch.
The total loss averages over all bifurcations in the training patch:
$\mathcal{L}_{\mathrm{softBCS}} = \frac{1}{|\mathcal{C}'|}\sum_{c \in \mathcal{C}'} \ell_c$.

\subsection{Base Model and Training Setup}

We use CT-FM~\cite{pai_ctfm_2025}, a SegResNet~\cite{myronenko_segresnet_2019} encoder (${\sim}$78M parameters) pretrained on 148K CT scans via self-supervised learning.
All experiments start from an ImageCAS-pretrained checkpoint that was fine-tuned with a multi-task loss combining DiceCE segmentation and Frangi vesselness prediction~\cite{frangi_multiscale_1998} (denoted $\mathcal{L}_v$, our ``L1'' baseline).
The composite training loss is:
\begin{equation}\label{eq:loss}
\mathcal{L} = \mathcal{L}_v + \lambda_s \, \mathcal{L}_{\mathrm{softBCS}},
\end{equation}
where $\lambda_s$ controls the topology loss weight.
Training uses ImageCAS~\cite{zeng_imagecas_2023} (638 train, 112 val, 250 test), $192^3$ patches with random pos/neg cropping, 30 epochs, and cosine-annealed learning rate.

\subsection{Adaptation Strategies}
\label{sec:strategies}

We evaluate five strategies spanning a parameter-efficiency spectrum (Table~\ref{tab:strategies}).
All start from the same L1 checkpoint, use identical data augmentation and training schedule, and differ only in which parameters are trainable.

\begin{table}[t]
\caption{Encoder adaptation strategies. All share the same 10M-parameter decoder.}
\label{tab:strategies}
\centering
\small
\begin{tabular}{llrr}
\toprule
Strategy & Encoder trainable params & Total trainable & LR \\
\midrule
Frozen (decoder-only)     & 0                        & 10.0M  & $10^{-5}$ \\
BitFit (norm params)      & ${\sim}$14K (BN bias+scale) & 10.0M  & $10^{-5}$ \\
Adapters (skip bottleneck)& ${\sim}$175K              & 10.2M  & $10^{-5}$ \\
Partial (last enc.\ block) & 17.7M                   & 27.7M  & $5{\times}10^{-5}$ \\
Full (all parameters)     & 77.8M                    & 87.8M  & $10^{-4}$ \\
\bottomrule
\end{tabular}
\end{table}

Each strategy is crossed with $\lambda_s \in \{0.03, 0.05, 0.10\}$.
BitFit and adapter experiments are run at $\lambda_s{=}0.05$ only; full fine-tuning at $\lambda_s{=}0.10$ only (it diverged at lower weights).
This yields 11 experimental cells plus the L1 baseline ($\lambda_s{=}0$).

\subsection{Representational Analysis}

To understand \emph{why} certain strategies preserve performance, we measure how much the encoder changes during fine-tuning.

\paragraph{CKA similarity~\cite{kornblith_cka_2019}.}
We compute linear Centered Kernel Alignment between the L1 baseline encoder and each fine-tuned encoder at five resolution levels (32ch/96$^3$ to 512ch/6$^3$).
CKA${\approx}1$ indicates preserved representations; low CKA indicates feature drift.

\paragraph{Linear probing.}
We train logistic regression classifiers on frozen encoder features to predict whether a voxel neighbourhood contains a bifurcation.
Higher probe accuracy at a given level indicates that bifurcation-relevant information is linearly accessible there.
We probe both encoder and decoder features.

\subsection{Cross-Domain Transfer: ASOCA}
\label{sec:asoca}

ASOCA~\cite{gharleghi_asoca_2022} provides 40 CTCA scans (20 normal, 20 diseased) with major-vessel-only annotations, contrasting ImageCAS's full coronary tree labels.
This annotation protocol mismatch means the pretrained model systematically includes side branches absent from ASOCA's ground truth.

\paragraph{Zero-shot evaluation.}
We apply all ImageCAS-trained models directly to ASOCA (threshold 0.5) to quantify the mismatch.

\paragraph{Low-shot fine-tuning.}
We fine-tune using $k \in \{1, 5, 10\}$ labelled ASOCA scans per group ($2k$ total), with topology-aware supervision ($\lambda_s{=}0.05$).
The encoder and first two decoder levels are frozen, adapting only the last two decoder levels---motivated by the insight that the model needs to suppress side branches rather than learn new features.
We sweep the binarisation threshold $t \in [0.3, 0.9]$ and report results at the optimal threshold per $k$.

\subsection{Evaluation Metrics}

We report Dice (volumetric overlap), clDice~\cite{shit_cldice_2021} (centreline overlap), BCS~\cite{anonymous_bcs_2026} (bifurcation connectedness, $\tau{=}3$, $l{=}8$), HD95 (95th-percentile Hausdorff distance, mm), and Betti$_0$ error ($|\beta_0(P) - \beta_0(G)|$, connected component mismatch).

%----------------------------------------------------------------------
\section{Results}

\subsection{ImageCAS Ablation: Adaptation Strategy $\times$ BCS Weight}

Table~\ref{tab:ablation} presents the full ablation grid on $n{=}250$ ImageCAS test cases.

\begin{table}[t]
\caption{Ablation grid on ImageCAS ($n{=}250$ test). Best per-column in \textbf{bold}; $\dagger$ marks significant improvement over L1 baseline (paired bootstrap, $p{<}0.05$). Full fine-tuning (EXP-045) catastrophically overfits despite achieving the highest BCS.}
\label{tab:ablation}
\centering
\small
\begin{tabular}{llccccc}
\toprule
Adaptation & $\lambda_s$ & Dice $\uparrow$ & clDice $\uparrow$ & BCS $\uparrow$ & HD95\,(mm) $\downarrow$ & Betti$_0$ $\downarrow$ \\
\midrule
\multicolumn{2}{l}{\textit{L1 baseline ($\lambda_s{=}0$)}} & 0.790 & 0.867 & 0.690 & 10.5 & 4.70 \\
\midrule
\multirow{3}{*}{Frozen}
 & 0.03 & \textbf{0.793}$^\dagger$ & 0.857 & 0.716$^\dagger$ & 8.7$^\dagger$ & 9.44 \\
 & 0.05 & 0.792$^\dagger$ & 0.856 & 0.722$^\dagger$ & 8.8$^\dagger$ & 9.11 \\
 & 0.10 & 0.785 & 0.844 & 0.730$^\dagger$ & 10.2 & 13.20 \\
\midrule
\multirow{3}{*}{Partial}
 & 0.03 & \textbf{0.793}$^\dagger$ & 0.858 & 0.715$^\dagger$ & 9.4$^\dagger$ & 8.44 \\
 & 0.05 & 0.792$^\dagger$ & \textbf{0.861} & 0.724$^\dagger$ & \textbf{8.4}$^\dagger$ & \textbf{8.13} \\
 & 0.10 & 0.789 & 0.858 & 0.714$^\dagger$ & 8.5$^\dagger$ & 7.95 \\
\midrule
BitFit & 0.05 & 0.787 & 0.851 & 0.720$^\dagger$ & 9.5$^\dagger$ & 10.93 \\
Adapters & 0.05 & 0.781 & 0.840 & 0.715$^\dagger$ & 13.6 & 14.44 \\
\midrule
Full & 0.10 & 0.748 & 0.778 & \textbf{0.739} & 33.6 & 19.85 \\
\bottomrule
\end{tabular}
\end{table}

\paragraph{Effect of BCS weight.}
Within both frozen and partial rows, increasing $\lambda_s$ monotonically raises BCS but progressively lowers Dice.
At $\lambda_s{=}0.03$, BCS improves by $+2.6$\,pp over L1 with no Dice penalty; at $\lambda_s{=}0.10$, BCS reaches its peak ($+4.0$\,pp) but Dice drops by $0.5$--$0.8$\,pp.
The sweet spot is $\lambda_s \in [0.03, 0.05]$.

\paragraph{Frozen vs.\ partial.}
At matched $\lambda_s$, frozen and partial strategies achieve nearly identical Dice (${\leq}0.1$\,pp difference).
Partial fine-tuning offers slightly better HD95 and Betti$_0$ at $\lambda_s{=}0.05$ (8.4\,mm and 8.1 vs.\ 8.8\,mm and 9.1), suggesting the unfrozen encoder block contributes to surface precision, but the overall benefit is marginal.

\paragraph{Full fine-tuning.}
Despite achieving the highest BCS (0.739), full fine-tuning catastrophically overfits: validation Dice reaches 0.807 but test Dice drops to 0.748 (6\,pp gap), and HD95 triples to 33.6\,mm.
This makes it unsuitable despite its apparent topology advantage.
Figure~\ref{fig:scatter} visualises the change in Dice and BCS relative to the L1 baseline for all nine experimental cells.
Decoder-only and partial strategies gain $+2.4$ to $+4.0$\,pp BCS with near-zero Dice change; BitFit and adapters achieve similar BCS gains ($+3.0$ and $+2.5$\,pp) but with small Dice penalties ($-0.3$ and $-0.9$\,pp); full fine-tuning gains the most BCS ($+4.9$\,pp) but loses $-4.2$\,pp Dice.

\begin{figure}[t]
\centering
\includegraphics[width=0.95\linewidth]{../figures/fig1_dice_bcs_scatter.pdf}
\caption{Change in Dice and BCS relative to L1 baseline across all nine experimental cells. Green bars: BCS gain; blue/purple bars: Dice change. All parameter-efficient strategies improve BCS ($+2.4$ to $+4.0$\,pp) with ${\leq}1$\,pp Dice cost. Full fine-tuning (red) gains the most BCS ($+4.9$\,pp) but catastrophically loses $-4.2$\,pp Dice. Decoder-only (frozen) achieves the best Dice--BCS balance.}
\label{fig:scatter}
\end{figure}

\subsection{CKA Similarity Analysis}

Table~\ref{tab:cka} shows CKA similarity between the L1 encoder and each fine-tuned encoder at five resolution levels.

\begin{table}[t]
\caption{CKA similarity to L1 baseline at each encoder level. Higher values indicate preserved representations. All parameter-efficient strategies maintain CKA${\geq}0.84$; full fine-tuning shows severe drift. Frozen CKA is invariant to $\lambda_s$ (encoder weights unchanged).}
\label{tab:cka}
\centering
\small
\begin{tabular}{lccccc}
\toprule
 & \multicolumn{5}{c}{Encoder Level (channels, spatial resolution)} \\
\cmidrule(lr){2-6}
Strategy & 32ch, 96$^3$ & 64ch, 48$^3$ & 128ch, 24$^3$ & 256ch, 12$^3$ & 512ch, 6$^3$ \\
\midrule
Frozen ($\lambda_s{=}0.05$)  & 1.00 & 0.99 & 0.99 & 0.98 & 0.88 \\
BitFit ($\lambda_s{=}0.05$)  & 1.00 & 0.99 & 0.99 & 0.98 & 0.89 \\
Adapter ($\lambda_s{=}0.05$) & 1.00 & 0.99 & 0.99 & 0.97 & \textbf{0.90} \\
Partial ($\lambda_s{=}0.05$) & 1.00 & 0.99 & 0.99 & 0.97 & 0.84 \\
Full ($\lambda_s{=}0.10$)    & 0.95 & 0.73 & 0.67 & 0.44 & 0.37 \\
\bottomrule
\end{tabular}
\end{table}

Frozen strategies produce identical CKA profiles regardless of $\lambda_s$ (expected, as the encoder is fixed).
The slight CKA drop at the bottleneck (0.88) reflects batch normalisation running statistics drifting during decoder training.
BitFit and adapter strategies achieve marginally \emph{higher} bottleneck CKA (0.89 and 0.90) than frozen, despite modifying encoder parameters---explicit optimisation of normalisation statistics or addition of lightweight modules appears to stabilise the representation.
Partial unfreezing allows more drift at the bottleneck (0.84) but preserves all other levels.
Full fine-tuning drastically reduces CKA at every level, with the bottleneck reaching 0.37---a near-complete reorganisation of the deepest features.
This confirms that the pretrained encoder already contains topology-relevant representations; all parameter-efficient strategies preserve them, while full fine-tuning \emph{destroys} them.
Figure~\ref{fig:cka} visualises the full CKA matrix across all nine experimental cells, revealing that frozen strategies are invariant to $\lambda_s$ (identical encoder weights), while the gradient of drift from partial to full is steep and concentrated at deeper levels.

\begin{figure}[t]
\centering
\includegraphics[width=0.85\linewidth]{../figures/fig1_cka_heatmap.pdf}
\caption{How much does fine-tuning change the encoder? The blue band shows the range across all eight parameter-efficient models (frozen, BitFit, adapter, partial); the red line shows full fine-tuning. Parameter-efficient strategies preserve encoder representations at all depths (CKA${>}0.82$); full fine-tuning progressively destroys them, reaching CKA${=}0.37$ at the deepest level.}
\label{fig:cka}
\end{figure}

\subsection{Linear Probing: Encoder and Decoder}

Table~\ref{tab:probe} reports bifurcation detection AUC from linear probes on frozen encoder and decoder features.

\begin{table}[t]
\caption{Linear probe AUC for bifurcation detection at each encoder and decoder level. \textbf{Bold}: highest AUC per model. All parameter-efficient strategies peak at encoder level~1; full fine-tuning shifts the peak to level~2.}
\label{tab:probe}
\centering
\small
\begin{tabular}{lccccc}
\toprule
 & \multicolumn{5}{c}{Encoder Level (AUC)} \\
\cmidrule(lr){2-6}
Model & 32ch, 96$^3$ & 64ch, 48$^3$ & 128ch, 24$^3$ & 256ch, 12$^3$ & 512ch, 6$^3$ \\
\midrule
L1 baseline             & 0.684 & \textbf{0.813} & 0.772 & 0.705 & 0.620 \\
Frozen ($\lambda_s{=}0.05$)  & 0.686 & \textbf{0.811} & 0.763 & 0.705 & 0.623 \\
BitFit ($\lambda_s{=}0.05$)  & 0.687 & \textbf{0.812} & 0.765 & 0.706 & 0.621 \\
Adapter ($\lambda_s{=}0.05$) & 0.687 & \textbf{0.811} & 0.763 & 0.705 & 0.614 \\
Partial ($\lambda_s{=}0.05$) & 0.687 & \textbf{0.811} & 0.763 & 0.700 & 0.602 \\
Full ($\lambda_s{=}0.10$)    & 0.680 & 0.773 & \textbf{0.794} & 0.694 & 0.638 \\
\midrule
 & \multicolumn{4}{c}{Decoder Level (AUC)} & \\
\cmidrule(lr){2-5}
 & 256ch, 12$^3$ & 128ch, 24$^3$ & 64ch, 48$^3$ & 32ch, 96$^3$ & \\
\midrule
L1 baseline             & 0.670 & 0.710 & \textbf{0.760} & 0.705 & \\
Frozen ($\lambda_s{=}0.05$)  & 0.667 & 0.710 & \textbf{0.756} & 0.690 & \\
BitFit ($\lambda_s{=}0.05$)  & 0.669 & 0.712 & \textbf{0.758} & 0.690 & \\
Adapter ($\lambda_s{=}0.05$) & 0.676 & 0.708 & \textbf{0.755} & 0.694 & \\
Partial ($\lambda_s{=}0.05$) & 0.669 & 0.708 & \textbf{0.750} & 0.692 & \\
Full ($\lambda_s{=}0.10$)    & 0.676 & \textbf{0.766} & 0.754 & 0.665 & \\
\bottomrule
\end{tabular}
\end{table}

In the encoder, all parameter-efficient strategies (frozen, BitFit, adapter, partial) peak at level~1 (64ch, 48$^3$; AUC${\approx}$0.81), nearly identical to the L1 baseline---confirming that these strategies preserve the encoder's bifurcation-discriminative features regardless of how many encoder parameters are trainable.
Full fine-tuning shifts the peak to level~2 (AUC~0.79), indicating feature hierarchy reorganisation rather than enhancement.

In the decoder, all models peak at level~2 (64ch, 48$^3$; AUC${\approx}$0.75--0.76).
Full fine-tuning again shows a shifted peak (level~1, AUC~0.77).
The overall similarity across decoder probes suggests that the Soft BCS loss refines how bifurcation information is routed to the segmentation output rather than dramatically relocating it within the network.

\subsection{ASOCA: Zero-Shot Transfer under Protocol Mismatch}
\label{sec:asoca_results}

Table~\ref{tab:zeroshot} reports zero-shot performance on ASOCA ($N{=}40$).

\begin{table}[t]
\caption{Zero-shot transfer to ASOCA ($N{=}40$). ImageCAS-pretrained models over-segment: high BCS (branches preserved) but low Dice (side branches penalised by major-vessel GT). Higher $\lambda_s$ increases BCS but further reduces Dice.}
\label{tab:zeroshot}
\centering
\small
\begin{tabular}{llccccc}
\toprule
Adaptation & $\lambda_s$ & Dice $\uparrow$ & BCS $\uparrow$ & HD95\,(mm) $\downarrow$ & Betti$_0$ $\downarrow$ \\
\midrule
\multicolumn{2}{l}{\textit{L1 baseline ($\lambda_s{=}0$)}} & \textbf{0.481} & 0.794 & \textbf{15.7} & \textbf{7.1} \\
\midrule
\multirow{3}{*}{Frozen}
 & 0.03 & 0.461 & 0.821 & 17.3 & 11.3 \\
 & 0.05 & 0.446 & 0.824 & 21.1 & 13.8 \\
 & 0.10 & 0.430 & \textbf{0.838} & 25.9 & 21.4 \\
\midrule
BitFit & 0.05 & 0.447 & 0.834 & 19.8 & 12.1 \\
\midrule
\multirow{3}{*}{Partial}
 & 0.03 & 0.453 & 0.815 & 22.6 & 14.5 \\
 & 0.05 & 0.431 & 0.829 & 21.9 & 14.8 \\
 & 0.10 & 0.446 & 0.830 & 20.4 & 14.4 \\
\midrule
Full & 0.10 & 0.423 & 0.831 & 43.0 & 21.6 \\
\bottomrule
\end{tabular}
\end{table}

All models achieve high BCS (${\geq}0.79$) on ASOCA despite low Dice (${\leq}0.48$), confirming systematic over-segmentation due to annotation protocol mismatch rather than image-domain shift: the models predict too \emph{much} rather than too little.
Higher $\lambda_s$ further increases BCS ($+4.4$\,pp for frozen $\lambda_s{=}0.10$) at the cost of Dice ($-5.1$\,pp), amplifying the mismatch.
A threshold sweep ($t \in [0.3, 0.9]$) reveals that BCS remains remarkably stable (${\pm}2$\,pp) while Dice improves monotonically with higher thresholds.
At $t{=}0.9$, L1 recovers $+18$\,pp Dice (0.479${\to}$0.665) with only $-1.6$\,pp BCS loss.
Importantly, L1 outperforms all BCS-trained models at every threshold on ASOCA---BCS training reinforces branch preservation, which is precisely the opposite of what major-vessel-only annotations require.
Figure~\ref{fig:threshold} shows the full threshold sweep: Dice increases monotonically with threshold while BCS remains stable (${\pm}2$\,pp), offering a cost-free calibration knob for protocol mismatch.

\begin{figure}[t]
\centering
\includegraphics[width=\linewidth]{../figures/fig3_asoca_threshold.pdf}
\caption{Threshold calibration on ASOCA: a free lunch for protocol mismatch. \textbf{(a)}~Raising the binarisation threshold from 0.5 to 0.9 recovers $+18.6$\,pp Dice (green shading marks the improvement zone); L1 outperforms all BCS-trained models at every threshold. \textbf{(b)}~BCS remains stable across thresholds (${\pm}2$\,pp), confirming that overlap gains come at no connectivity cost.}
\label{fig:threshold}
\end{figure}

\subsection{ASOCA: Low-Shot Adaptation}

We fine-tune the L1 decoder on ASOCA with $k \in \{1, 5, 10\}$ scans per group (total $2k$ training cases), freezing the encoder and first two decoder levels (only the last two decoder levels are trainable).
This shallow adaptation follows the insight that the model needs to \emph{suppress} side branches rather than learn new features.
Table~\ref{tab:lowshot} reports topology-aware fine-tuning ($\lambda_s{=}0.05$) at the optimal binarisation threshold per $k$.

\begin{table}[t]
\caption{Low-shot fine-tuning on ASOCA with topology-aware supervision ($\lambda_s{=}0.05$). Only the last two decoder levels are trainable. Results at the optimal binarisation threshold per $k$. $k{=}1$ achieves the best Dice--BCS balance; more data overfits to ASOCA's sparse annotations.}
\label{tab:lowshot}
\centering
\small
\begin{tabular}{ccccccc}
\toprule
$k$ & $t$ & $n_{\mathrm{eval}}$ & Dice $\uparrow$ & BCS $\uparrow$ & HD95\,(mm) $\downarrow$ & Betti$_0$ $\downarrow$ \\
\midrule
\multicolumn{2}{l}{\textit{Zero-shot ($t{=}0.5$)}} & 40 & 0.481 & 0.794 & 15.7 & 7.1 \\
\multicolumn{2}{l}{\textit{Zero-shot ($t{=}0.9$)}} & 40 & 0.665 & 0.791 & 13.0 & 10.0 \\
\midrule
1  & 0.7 & 34 & \textbf{0.706} & \textbf{0.752} & \textbf{10.6} & 13.1 \\
5  & 0.4 & 26 & 0.695 & 0.699 & 26.8 & --- \\
10 & 0.4 & 16 & 0.671 & 0.751 & 46.4 & 64.6 \\
\bottomrule
\end{tabular}
\end{table}

Remarkably, $k{=}1$ (just two training scans) at $t{=}0.7$ achieves the best combined performance: Dice~0.706 ($+22.5$\,pp over zero-shot), BCS~0.752, and HD95~10.6\,mm.
Adding more ASOCA training data ($k{=}5$, $k{=}10$) paradoxically degrades performance---BCS drops by $5$--$5.3$\,pp and HD95 increases sharply---indicating that the model overfits to ASOCA's sparser annotation protocol, destroying the pretrained topology-aware features.
The shallow decoder freeze (depth~2) is critical: our earlier experiment with full decoder fine-tuning (7 training cases, $\lambda_s{=}0.30$) achieved similar Dice (0.701) but catastrophically lower BCS (0.679) and Betti error (36.4 vs.\ 13.1).

%----------------------------------------------------------------------
\section{Discussion}

\paragraph{Decoder-only adaptation is sufficient.}
The central finding of the ImageCAS ablation is that topology-aware fine-tuning does not require modifying the pretrained encoder.
Decoder-only fine-tuning achieves the best balance (Dice~0.793, BCS~0.716); BitFit and adapter strategies achieve comparable BCS ($+3.0$ and $+2.5$\,pp) but at a Dice cost ($-0.3$ and $-0.9$\,pp respectively), while full fine-tuning catastrophically overfits.
CKA analysis provides a mechanistic explanation: the pretrained encoder already encodes topology-relevant features (linear probes detect bifurcations at AUC~0.81 at all parameter-efficient strategies), and the decoder simply learns to exploit them for connectivity.
Notably, adapter modules achieve the highest bottleneck CKA (0.90), yet this superior representation preservation does not translate to measurably better segmentation---confirming that the encoder is already sufficient and the decoder is where adaptation matters.
Full fine-tuning destroys these representations (CKA drops to 0.37 at the bottleneck), explaining the 6\,pp val--test Dice gap.
This mirrors findings from NLP and 2D transfer learning~\cite{kumar_finetuning_2022}, extending them to 3D topology-aware medical image segmentation.

\paragraph{BCS weight controls the overlap--connectivity trade-off.}
Within the effective strategies (frozen and partial), increasing $\lambda_s$ monotonically improves BCS at the cost of Dice.
The effect is modest: from $\lambda_s{=}0.03$ to $\lambda_s{=}0.10$, BCS improves by $+1.4$\,pp but Dice drops by $0.8$\,pp.
This suggests a Pareto frontier that practitioners can navigate based on downstream requirements---higher $\lambda_s$ for connectivity-critical applications (FFR-CT, stent planning), lower $\lambda_s$ for volumetric accuracy.

\paragraph{Annotation protocol mismatch.}
Zero-shot transfer to ASOCA reveals a striking dissociation: BCS is high (0.794) but Dice is low (0.481).
The model correctly identifies vascular branches, including those that are part of the full coronary tree but absent from ASOCA's major-vessel annotations.
Threshold tuning ($t{=}0.9$) recovers $+18$\,pp Dice with only $-1.6$\,pp BCS loss, while $k{=}1$ low-shot fine-tuning with topology supervision achieves Dice~0.706 at BCS~0.752.
Crucially, more training data ($k{=}5$, $k{=}10$) \emph{degrades} performance, suggesting that the model overfits to the sparser annotation protocol.
Furthermore, starting low-shot adaptation from a BCS-trained checkpoint (frozen $\lambda_s{=}0.05$) rather than the L1 baseline yields worse results (Dice~0.680, BCS~0.624 vs.\ 0.706/0.752), indicating that topology-aware pretraining on full-tree annotations actively hinders adaptation to a major-vessel protocol---the learned branch-preservation features resist suppression.
This highlights that annotation protocol mismatch requires careful calibration rather than brute-force supervised adaptation.

\paragraph{Limitations.}
This study evaluates a single foundation model (CT-FM SegResNet); generality to other architectures remains to be tested.
ASOCA's small size ($N{=}40$) limits statistical power for the low-shot analysis, and the $k{=}1$ result uses only 34 evaluation cases.
The Soft BCS loss requires precomputed skeleton labels, which may not always be available.

%----------------------------------------------------------------------
\section{Conclusion}

We presented a systematic study of encoder adaptation strategies for topology-aware fine-tuning of a CT foundation model for coronary artery segmentation.
Decoder-only fine-tuning achieves the best overlap--connectivity balance, preserving pretrained encoder features (CKA${\geq}0.84$) while enabling the decoder to learn bifurcation-level connectivity.
BitFit and adapter strategies confirm that even minimal encoder modification yields equivalent representations---the encoder already contains topology-relevant features that the decoder need only learn to exploit.
Full fine-tuning catastrophically overfits despite achieving the highest BCS.
On ASOCA, annotation protocol mismatch causes systematic over-segmentation that BCS training amplifies rather than mitigates; minimal adaptation ($k{=}1$ with shallow decoder fine-tuning and threshold calibration) recovers Dice from 0.481 to 0.706, while starting from a BCS-pretrained checkpoint or adding more training data paradoxically degrades performance.
These findings suggest that practitioners should freeze the encoder and focus adaptation on the decoder when adding topology-aware objectives, and that cross-domain deployment under annotation protocol mismatch requires calibration rather than additional supervised training.

%----------------------------------------------------------------------
\begin{credits}
\subsubsection{\ackname}
This research was supported by [funder and project code to be disclosed upon publication].

\subsubsection{\discintname}
The authors have no competing interests to declare that are relevant to the content of this article.
\end{credits}

\bibliographystyle{splncs04}
\bibliography{references}

\end{document}
